\chapter{Lung Disease}

\section*{Patient focus}
Chest imaging is one of the most common types of scans. A chest X-ray can show pneumonia, fluid around the lungs, or a collapsed lung. A CT scan gives much more detail about the lung tissue itself and can identify scarring (fibrosis), emphysema (damaged air sacs), or small nodules that might need follow-up. If you smoke or have a family history of lung cancer, screening CT may be recommended. The way your lungs look on a scan helps your doctor decide on treatment and monitor your condition over time.


\label{ch:ch62_lung_disease}

\section{Imaging}
Radiography is useful for many acute and follow-up questions; CT characterizes parenchyma, airways, nodules, vessels, and pleura; MRI and nuclear methods have selected roles. Use high-resolution or angiographic protocols only for defined indications.

\section{COPD Imaging}
\begin{itemize}
  \item \textbf{CXR}: hyperinflation (flattened diaphragms, increased retrosternal space), bullae, pulmonary artery enlargement
  \item \textbf{CT}: emphysema quantification (centrilobular vs panlobular vs paraseptal), airway wall thickening, bronchiectasis, bullous disease
  \item \textbf{CT for lung volume reduction surgery (LVRS)}: map emphysema distribution (upper vs lower lobe predominance) to determine surgical candidacy
  \item \textbf{CT densitometry}: quantitative assessment of emphysema extent---used in endobronchial valve planning
\end{itemize}

\section{Interstitial Lung Disease Patterns}

High-resolution CT (HRCT) is the cornerstone for ILD diagnosis. Pattern recognition combined with clinical context allows non-diagnosis in many cases.

\subsection{Usual Interstitial Pneumonia (UIP)}
\begin{itemize}
  \item \textbf{Distribution}: basal and peripheral predominance
  \item \textbf{Features}: reticular opacities, honeycombing (clustered cystic air spaces), traction bronchiectasis, minimal ground-glass
  \item \textbf{Temporal heterogeneity}: areas of fibrosis adjacent to relatively normal lung
  \item \textbf{Most common cause}: idiopathic pulmonary fibrosis (IPF)
\end{itemize}

\subsection{Non-Specific Interstitial Pneumonia (NSIP)}
\begin{itemize}
  \item \textbf{Distribution}: basal and peripheral, but more uniform than UIP
  \item \textbf{Features}: ground-glass opacities (cellular NSIP) or fine reticulation with traction bronchiectasis (fibrotic NSIP), \textbf{no honeycombing}
  \item \textbf{Most common causes}: connective tissue disease, drug toxicity, hypersensitivity pneumonitis
\end{itemize}

\subsection{Organizing Pneumonia (OP)}
\begin{itemize}
  \item \textbf{Distribution}: patchy, subpleural or peribronchovascular
  \item \textbf{Features}: consolidation with air bronchograms, reversed halo sign (atoll sign), migrating opacities
  \item \textbf{Most common cause}: cryptogenic organizing pneumonia (COP); also secondary to infection, drugs, autoimmune disease
\end{itemize}

\begin{pearlbox}[UIP vs NSIP on HRCT]
Honeycombing is the key differentiator. UIP shows basal-predominant honeycombing with temporal heterogeneity; NSIP shows ground-glass with fine reticulation and traction bronchiectasis but \textbf{no honeycombing}. If honeycombing is present, IPF is most likely until proven otherwise. A multidisciplinary discussion (MDD) is recommended when HRCT is indeterminate.
\end{pearlbox}

\begin{table}[h]
\centering
\caption{HRCT patterns in ILD}
\begin{tabular}{|p{0.13\textwidth}|p{0.23\textwidth}|p{0.33\textwidth}|p{0.15\textwidth}|}
\hline
\textbf{Pattern} & \textbf{Distribution} & \textbf{Key features} & \textbf{Common cause} \\
\hline
HP & Upper/mid, peribronchovascular & Mosaic attenuation, air trapping, GGO & Hypersensitivity \\
LIP & Basal, peribronchovascular & GGO, cysts, nodules & HIV, Sjogren \\
NSIP & Basal, peripheral & GGO, fine reticulation, no honeycombing & CTD, drugs \\
OP & Patchy, subpleural & Consolidation, reversed halo sign & COP \\
UIP & Basal, peripheral & Honeycombing, reticulation, traction bronchiectasis & IPF \\
\hline
\end{tabular}
\end{table}

\section{Pulmonary Fibrosis}
\begin{itemize}
  \item \textbf{Progression indicators}: increasing honeycombing, traction bronchiectasis, volume loss, worsening reticulation
  \item \textbf{Drug-related fibrosis}: amiodarone (high attenuation), methotrexate, bleomycin, nitrofurantoin
  \item \textbf{Asbestos-related}: pleural plaques with lower lobe fibrosis (asbestosis); diffuse pleural thickening
\end{itemize}

\section{Lung Cancer Screening}

\subsection{NLST Criteria}
The National Lung Screening Trial demonstrated 20\% mortality reduction with annual low-dose CT (LDCT):
\begin{itemize}
  \item Age 55--74
  \item $\geq$30 pack-year smoking history
  \item Currently smoking or quit within 15 years
\end{itemize}

\subsection{Lung-RADS}
The Lung Imaging Reporting and Data System standardizes screening CT reporting:

\begin{table}[h]
\centering
\caption{Lung-RADS categories}
\begin{tabular}{|l|l|l|}
\hline
\textbf{Category} & \textbf{Description} & \textbf{Action} \\
\hline
1 & Negative, no nodules & Annual screening \\
2 & Benign appearance (solid $<$6 mm, GGO $<$30 mm) & Annual screening \\
3 & Probably benign (solid 6--$<$8 mm, GGO $\geq$30 mm) & 6-month follow-up \\
4A & Suspicious (solid 8--$<$15 mm) & 3-month follow-up or PET \\
4B & Highly suspicious (solid $\geq$15 mm, spiculated) & Tissue diagnosis \\
4X & Additional features & Further work-up \\
\hline
\end{tabular}
\end{table}

\begin{pearlbox}[Screening vs diagnostic CT]
Screening LDCT uses lower radiation dose (1--2 mSv) and no IV contrast. If a significant abnormality is found, the patient transitions from screening to diagnostic work-up with higher-dose CT, contrast, and potentially PET/CT or biopsy. A screening CT is not a substitute for diagnostic evaluation.
\end{pearlbox}

\section{Pleural Disease}
\begin{itemize}
  \item \textbf{Pleural effusion}: CXR (blunting of costophrenic angle, meniscus sign); CT characterises loculation and pleural thickening; ultrasound for procedural guidance
  \item \textbf{Pneumothorax}: CXR (visceral pleural line, absent lung markings); CT more sensitive (anterior in supine patients); E-FAST for occult pneumothorax
  \item \textbf{Pleural thickening}: diffuse vs focal; asbestos-related pleural disease vs mesothelioma vs metastatic disease
  \item \textbf{Empyema}: enhancing pleural rind, loculated fluid, split pleura sign on CT---may need surgical decortication
\end{itemize}

\section{Patterns}
Describe consolidation, ground glass, nodules, reticulation, honeycombing, airway disease, emphysema, edema, pleural disease, and distribution. Correlate with symptoms, exposure, immune status, pulmonary function, and prior imaging.

\section{Communication}
Report pneumothorax, pulmonary embolism, severe infection, airway obstruction, major hemoptysis source, or rapidly progressive disease promptly. Recommendations should follow current guideline pathways and patient risk.

\begin{pearlbox}[Rapidly progressive ILD]
If HRCT shows subacute worsening of fibrosis with new ground-glass and clinical deterioration, consider acute exacerbation of IPF, drug toxicity, or superimposed infection. Communicate urgently---mortality from acute exacerbation of IPF exceeds 50\%.
\end{pearlbox}

\clearpage
